\documentclass[11pt,]{article}
\usepackage[margin=1in]{geometry}
\newcommand*{\authorfont}{\fontfamily{phv}\selectfont}
\usepackage[]{mathpazo}
\usepackage{multicol}
\usepackage{abstract}
\usepackage{array}
\usepackage{ragged2e}  % For ragged formatting
\renewcommand{\abstractname}{}    % clear the title
\renewcommand{\absnamepos}{empty} % originally center
\newcommand{\blankline}{\quad\pagebreak[2]}

\providecommand{\tightlist}{%
  \setlength{\itemsep}{0pt}\setlength{\parskip}{0pt}} 
\usepackage{longtable,booktabs}
      
\usepackage{parskip}
\usepackage{titlesec}
\titlespacing\section{0pt}{12pt plus 4pt minus 2pt}{6pt plus 2pt minus 2pt}
\titlespacing\subsection{0pt}{12pt plus 4pt minus 2pt}{6pt plus 2pt minus 2pt}

\titleformat*{\subsubsection}{\normalsize\itshape}

\usepackage{titling}
\setlength{\droptitle}{-.15cm}

%\setlength{\parindent}{0pt}
%\setlength{\parskip}{6pt plus 2pt minus 1pt}
%\setlength{\emergencystretch}{3em}  % prevent overfull lines 

\usepackage[T1]{fontenc}
\usepackage[utf8]{inputenc}

\usepackage{fancyhdr}
\pagestyle{fancy}
\usepackage{lastpage}
\renewcommand{\headrulewidth}{0.3pt}
\renewcommand{\footrulewidth}{0.0pt} 
\lhead{}
\chead{}
\rhead{\footnotesize GY306(W): Sedimentology and Stratigraphy -- Fall
2026}
\lfoot{}
\cfoot{\small \thepage /\pageref*{LastPage}}
\rfoot{}

\fancypagestyle{firststyle}
{
\renewcommand{\headrulewidth}{0pt}%
   \fancyhf{}
   \fancyfoot[C]{\small \thepage/\pageref*{LastPage}}
}

%\def\labelitemi{--}
%\usepackage{enumitem}
%\setitemize[0]{leftmargin=25pt}
%\setenumerate[0]{leftmargin=25pt}




\makeatletter
\@ifpackageloaded{hyperref}{}{%
\ifxetex
  \usepackage[setpagesize=false, % page size defined by xetex
              unicode=false, % unicode breaks when used with xetex
              xetex]{hyperref}
\else
  \usepackage[unicode=true]{hyperref}
\fi
}
\@ifpackageloaded{color}{
    \PassOptionsToPackage{usenames,dvipsnames}{color}
}{%
    \usepackage[usenames,dvipsnames]{color}
}
\makeatother
\hypersetup{breaklinks=true,
            bookmarks=true,
            pdfauthor={},
             pdfkeywords = {},  
            pdftitle={GY306(W): Sedimentology and Stratigraphy},
            colorlinks=true,
            citecolor=blue,
            urlcolor=blue,
            linkcolor=magenta,
            pdfborder={0 0 0}}
\urlstyle{same}  % don't use monospace font for urls


\setcounter{secnumdepth}{0}

\usepackage{longtable}




\usepackage{setspace}

\title{GY306(W): Sedimentology and Stratigraphy}
\date{Fall 2026}


\begin{document}  

		\maketitle
		
	
		\thispagestyle{firststyle}

%	\thispagestyle{empty}
\textbf{\underline{Instructor}}
\begin{multicols}{2}

  \textbf{Dr.~Benjamin J. Linzmeier}\\
  Email: \href{mailto:blinzmeier@southalabama.edu}{\nolinkurl{blinzmeier@southalabama.edu}}\\
  Preferred contact: Canvas LMS messages\\
  Office Hours: TBD\\
  Office: Life Science Building room 344\\
  Correspondence policy: Within 24 hours on weekdays; 48 hours on
weekends\\
    \columnbreak
    
  \end{multicols}
	
\noindent \begin{tabular*}{\textwidth}{ @{\extracolsep{\fill}} lr @{\extracolsep{\fill}}}
\textbf{\underline{Class information}}\\
  Classroom: LCSB 337\\
  Class Hours: 12:20 - 1:10 pm MWF\\
    Lab room: LCSB 337\\
  Lab Hours: 1:25 pm - 3:55 pm W\\
    \\
	\end{tabular*}\\


\vspace{2mm}


\section{Course description}\label{course-description}

Why do Sedimentology and Stratigraphy (GY306{[}W{]}) matter? Sediments
house the record of climate at Earth's surface, influence climate
through feedbacks, and house important resources. In this course you
will learn how sediments are formed, moved, and preserved in the rock
record. We will discuss the way that paleoenvironments are interpreted
from observations of rocks on multiple scales. In addition, we will work
together to improve your writing skills with a focus on a few of the
genres of writing you may be expected to do as professional geologists.

\section{Writing requirements}\label{writing-requirements}

For this class to fulfill the (W) requirements of USA, you will be doing
a semester-long research project in an area related to sedimentology. To
build your skills in discipline-specific writing, we will be doing
smaller writing assignments throughout the semester. The final projects
will be peer-reviewed, revised, and then submitted for final grading and
assessment.

\section{Important Dates}\label{important-dates}

August 26th - Last day to drop without WD\\
October 31st - Last day to drop or withdraw. See the
\href{https://www.southalabama.edu/academiccalendar/}{Academic Calendar}
for more information.\\

\newpage

\section{Course goals}\label{course-goals}

\begin{enumerate}
\def\labelenumi{\arabic{enumi}.}
\item
  Observe, describe, classify, and communicate about sediments and
  sedimentary rocks.
\item
  Develop a framework for building hypotheses about the distributions of
  sediments, sedimentary structures, and sedimentary environments.
\item
  Communicate sedimentary geology to audiences using different genres of
  writing.
\item
  Read and summarize scientific papers focused on sedimentary geology
  and related fields.
\end{enumerate}

\section{Class textbooks}\label{class-textbooks}

We will have two books that you should have access to for the class this
semester. The book on writing below will be referenced and discussed
often. We will also refer to the general sedimentology and stratigraphy
book.

\subsubsection{For writing}\label{for-writing}

Heard, Stephen B., \textbf{The Scientist's Guide to Writing, 2nd
Edition: How to Write More Easily and Effectively throughout Your
Scientific Career.} 2nd edition. Princeton: Princeton University Press,
2022.

\subsubsection{For general Sedimentology and
Stratigraphy}\label{for-general-sedimentology-and-stratigraphy}

Gary Nichols, \textbf{Sedimentology and Stratigraphy}, 2nd edition,
Wiley-Blackwell, ISBN: 978-1-405-13592-4

or

Sam Boggs Jr., \textbf{Principles of Sedimentology and Stratigraphy},
5th edition, Pearson, ISBN: 0321643186

\subsection{Optional books}\label{optional-books}

\subsubsection{Field focused rock ID}\label{field-focused-rock-id}

Dorrik A.V. Stow, \textbf{Sedimentary Rocks in the Field: A Colour
Guide}, Manson Publishing Ltd, ISBN 10: 1874545693

\subsubsection{Improve the mechanics of your
writing}\label{improve-the-mechanics-of-your-writing}

William Strunk JR. and E.B. White, \textbf{The Elements of Style}, ISBN:
1594200696

\subsubsection{Lectures and papers}\label{lectures-and-papers}

Both lecture slides and papers will be available via Canvas. Slides will
be posted after lectures and assigned papers will be available in
Canvas.

\newpage

\begin{table}[!h]
\centering
\begin{tabular}{l>{\raggedright\arraybackslash}p{4cm}>{\raggedright\arraybackslash}p{4cm}>{\raggedright\arraybackslash}p{4cm}}
\toprule
Date & Lecture & Lab & Reading\\
\midrule
2026-08-19 & Introduction & Field & \\
2026-08-21 & Project and Writing &  & Part 1\\
2026-08-24 & Weathering &  & N.6 or B.1\\
2026-08-26 & Erosion & Sand description & \\
2026-08-28 & Writing behavior &  & Ch.3\&4\\
2026-08-31 & Rock cycle and clastics & Clastic ID & B.3 or N.2+Digital\\
2026-09-02 & Abstracts &  & Ch.9\\
2026-09-04 & Fluid Flow Dynamics &  & B.2 or N.4\\
2026-09-09 & Fluid Flow Dynamics & Flume & \\
2026-09-11 & Kinds of reading &  & Ch.28\\
2026-09-14 & Gravity flows &  & B.2.4 or N4.2/.3\\
2026-09-16 & Sed. Structures & Dauphin Island & \\
2026-09-18 & Tell your story &  & Ch.7\\
2026-09-21 & Continental deposits &  & B.8 or N.7-10\\
2026-09-23 & Continental deposits & Carbonate ID & \\
2026-09-25 & Parts to a paper &  & Ch.8\\
2026-09-28 & Marine deposits &  & B.10-11 or N.14-16\\
2026-09-30 & Marine deposits & Petrography & \\
2026-10-02 & Revision &  & Ch.23\\
2026-10-05 & Midterm review &  & \\
2026-10-07 & Midterm & No lab & \\
2026-10-12 & Relative Dating &  & B.12 or N.19\\
2026-10-14 & Facies & Stedman's Landing & \\
2026-10-16 & NSF Proposals &  & NSF Proposals\\
2026-10-19 & Presentations &  & \\
2026-10-21 & Presentations & Project figure & \\
2026-10-23 & Conferences &  & TBD\\
2026-10-26 & Radiometric dating &  & B.15 or N.21\\
2026-10-28 & Strat columns & Baldwin Outcrop & \\
2026-10-30 & Specious barriers &  & Specious barriers\\
2026-11-02 & Base level and strat &  & Wheeler 1964\\
2026-11-04 & Sed Basins & QGIS Map & B.16 or N.24\\
2026-11-06 & Diverse writing &  & Ch.26\\
2026-11-09 & Sequence Strat &  & B.13 or N.23\\
2026-11-11 & Sequence Strat & Correlation & \\
2026-11-13 & Draft discussion &  & \\
2026-11-16 & Correlation &  & B.12 or N.19\\
2026-11-18 & Subsurface strat &  & B.13 or N.22\\
2026-11-20 & Subsurface strat & Log interpretation & \\
2026-11-23 & Audiences &  & Ch.24\\
2026-11-30 & Final review &  & \\
2026-12-02 & Presentations & Presentations & \\
2026-12-04 & Lessons learned &  & Part VII\\
\bottomrule
\end{tabular}
\end{table}

\newpage

\section{Student conduct}\label{student-conduct}

\subsection{Attendance policy}\label{attendance-policy}

Students who complete less than 50\% of assignments/exams by number of
assignment and not weight, will be assigned a failing (F) grade.

If you stop attending or fail to

Review the `Attendance and Absences Policy', in the Undergraduate and
Graduate Bulletin for attendance and absences policy

An individual student is responsible for attending the classes in which
the student is officially enrolled. The quality of work will ordinarily
suffer from excessive absences.

For excessive absences (two or three consecutive class meetings) due to
illness, death in the family, or family emergency, the Dean of Students'
office should be advised. Absence notices will be sent to each
instructor notifying him of the reason for and the approximate length of
the absence. This notification does not constitute an excused absence.

Students receiving veterans' benefits are required to attend classes
according to the regulations of the Veterans Administration.

All international students on F-1 visas must comply with attendance
regulations as dictated by the Department of Justice, Immigration and
Naturalization Services. They must remain students in good standing with
at least twelve (12) hours per semester.

Students attending authorized off-campus functions or required
activities shall be excused by the responsible University official
through the Office of Academic Affairs. In case of doubt, instructors
may consult these lists in that office. Work missed as a result of these
excused absences may be made up.

\subsection{Grading}\label{grading}

\subsubsection{Late policy}\label{late-policy}

I will be grading assignments in batches to provide fair assessment to
everyone. This means late work can be disruptive. The penalty for late
work will be 25\% a day after the first 24 hours late. If you have a
reasonable excuse (illness, etc.) the late policy may be relaxed.

\subsubsection{Extra credit}\label{extra-credit}

Extra credit may be made available throughout the semester at my
discretion.

Assignments will be returned within a week of submission.

\begin{longtable}[]{@{}lr@{}}
\toprule\noalign{}
Grade & Range \\
\midrule\noalign{}
\endhead
\bottomrule\noalign{}
\endlastfoot
A & 90-100 \\
B & 80-89 \\
C & 70-79 \\
D & 60-69 \\
F & \textless{} 60 \\
F* & \textless{} 60 attendance \\
\end{longtable}

Assignment weighting will follow:

\begin{longtable}[]{@{}lr@{}}
\toprule\noalign{}
Item & Weight \\
\midrule\noalign{}
\endhead
\bottomrule\noalign{}
\endlastfoot
Writing scaffolding & 10 \% \\
Participation & 10 \% \\
Semester Project & 35 \% \\
Labratory Assignments & 25 \% \\
Midterm Exam & 10 \% \\
Final Exam & 10 \% \\
\end{longtable}

\newpage

\section{Writing tasks}\label{writing-tasks}

\subsection{Scaffolding tasks}\label{scaffolding-tasks}

Scaffolding writing tasks assigned to give you some practice and general
feedback to improve your writing.

\subsubsection{Abstract 1}\label{abstract-1}

First attempt at writing an abstract using the Nature paragraph
template.

\subsubsection{Figure Caption}\label{figure-caption}

Write the figure captions for several example figures using some
literature examples. Learn to describe figures in text.

\subsubsection{Abstract 2}\label{abstract-2}

Write the abstract for a second paper using the feedback from the first
abstract to improve your process.

\subsubsection{Peer Review}\label{peer-review}

Peer review of preprint or publication using example template. The goal
of this assignment is to focus on critiquing the logic put forward in
the paper with suggestions for refining it.

\subsection{Semester project}\label{semester-project}

Semester project outline of writing products.

\subsubsection{Project Story Summary}\label{project-story-summary}

Short outline of proposed project that answers the following questions.
(Based on S.B. Heard's book).\\
1) What is the central question?\\
2) Why is this question important?\\
3) What data (variables) are needed to answer this question?\\
4) What methods are used to get those data?\\
5) What analysis must be applied for the data to answer this question?\\
6) What results are likely?\\
7) How can results answer the question?\\
8) What does the answer tell us about the broader field?

\subsubsection{Proposal introduction
draft}\label{proposal-introduction-draft}

Carefully explain the knowledge that lead you to develop this hypothesis
and clearly outline the hypothesis to be tested. Consider these
questions: - What do we know about the question you are interested in
studying?\\
- What background knowledge do people need to assess if your hypothesis
can be tested?\\
- How could auxiliary hypotheses be tested?\\
- What specific data are needed to test the hypothesis and is there a
clear distinction between the hypothesis and alternatives?

\subsubsection{Proposal methods draft}\label{proposal-methods-draft}

This section outlines the required samples, locations, and data that
must be collected to test your hypothesis. Consider these questions:\\
- Where will samples or field observations come from?\\
- Will data be collected in the field or in a lab?\\
- What instruments are needed to perform the measurements?\\
- What difficulties may come from the interpretations of the
measurements?\\
- How many analyses are needed?

\subsubsection{Proposal intellectual merit
draft}\label{proposal-intellectual-merit-draft}

Create a section that specifically states how this proposal contributes
to the field of research and is potentially transformative (although
this is not as important for the assignment.) Consider the following
questions:\\
- What will the potential next steps be beyond testing this
hypothesis?\\
- How could the findings of this work have impacts on other areas of
research?

\subsubsection{Proposal broader impacts
draft}\label{proposal-broader-impacts-draft}

This section summarizes how your work will benefit society. This could
be through applications of your discovery, training of individuals, or
communication to people outside of the scientific community.\\
- Does the research have an economic or societal benefit?\\
- Does the work require the training of individuals?\\
- How would this work benefit the university or community?

\subsubsection{Project first draft}\label{project-first-draft}

Full proposal draft with figures. The important aspect of this
submission is assessment of the research plan and justification.

\subsubsection{Project peer review}\label{project-peer-review}

Review of a peer's proposal with the goal of improving the work by
increasing clarity, readability and communication.

\subsubsection{Project final draft}\label{project-final-draft}

Full, final project proposal with figures, time line and clear plan for
work. Focus is on clearly communicating to your peers that the work is
worth doing and that you can do it.

\subsubsection{Project presentation}\label{project-presentation}

Presentation of a peer's research proposal with a focus on justifying
funding their project. The goal is to convince the audience that the
project can be done and will have an important impact on society.

\section{On writing}\label{on-writing}

\subsection{Revising written work}\label{revising-written-work}

Nothing is perfect on the first draft. We all revise and rework our
written creations for a variety of reasons. We may not communicate the
point we want to get across to all readers. We may have forgotten or not
known important pieces of information initially. The tone of the writing
may be inappropriate for the audience. Word choice might obscure
meaning. The only way we can make our writing better is by revising it
based on constructive feedback from others or ourselves.

\subsection{Genres of Science writing}\label{genres-of-science-writing}

\subsubsection{Paper reviews}\label{paper-reviews}

Critical but constructive critique of unpublished research papers.
Informs the journal editor of how the work contributes to existing
literature and provides guide to authors and the editor for improvement
to the manuscript. Can be signed but often anonymous.

\subsubsection{Descriptive Reports}\label{descriptive-reports}

Descriptive reports are often made for industry or geological surveys.
These describe the composition and structure of rocks in an area.
Sometimes they include estimations of resource potentials or
environmental impacts of industry.

\subsubsection{Research papers}\label{research-papers}

These pieces of science writing typically test one or more closely
related hypotheses in the context of existing scientific knowledge. Most
have structure that is pre-defined including: abstract, introduction,
methods, results, discussion, and conclusion. Figures are an integral
part of most papers in the Geosciences and range from images of outcrops
to abstracted summaries of observations to plots of measurements.

\subsubsection{Conference abstracts}\label{conference-abstracts}

These are generally a step on the path to a full research publication
where results, interpretations, and context are publicly presented for
the first time.

\subsubsection{Grant Applications}\label{grant-applications}

Scientific work is generally funded with public or private funds that
are obtained by writing grants. Unlike research papers, grants must
persuade reviewers that the research planned is important, can be done
by the applicant, and can be achieved in the time frame of the funding
request.

\subsubsection{Science Communication}\label{science-communication}

Scientists also write general summaries of their work for non-specialist
audiences. These can be blog posts, or lesson plans that are shared for
educators. Some folks also write books for general audiences describing
the processes of science.

\subsection{Common processes}\label{common-processes}

For writing in the discipline, we focus on several specific questions to
guide our composition: 1) Who is our audience? 2) What is the one or two
sentence take-home point? 3) How quickly are we communicating the
information?

Science writing, like all writing, is socially situated and caries with
it the beliefs, values, and ideologies of the particular community and
culture.

\subsection{Tips and tricks for good
writing}\label{tips-and-tricks-for-good-writing}

Use a citation management software to organize your papers and create
bibiliographies. My preference is \href{https://www.zotero.org/}{Zotero}
because it is free, open source, and works across multiple operating
systems.

Write and then revise. Ask for feedback from peers or wait until you
have some distance from your writing and then revise.

Prompt yourself with questions that have short, concrete answers.

\subsection{Academic disruption
policy}\label{academic-disruption-policy}

The University of South Alabama's policy regarding Academic Disruption
is found in The Lowdown, the student handbook.
\url{http://www.southalabama.edu/lowdown/academicdisruption.shtml}.
Disruptive academic behavior is defined as individual or group conduct
that interrupts or interferes with any educational activity or
environment, infringes upon the rights and privileges of others, results
in or threatens the destruction of property and/or is otherwise
prejudicial to the maintenance of order in an academic environment. At
all times students will be cordial, courteous and respectful of faculty
members and fellow students. Cell phones, videotaping, and other
electronic devices are not allowed; however, you may use a laptop for
note taking. If your laptop is used for other purposes, the instructor
holds the right to revoke laptop use.

\subsection{Academic honesty}\label{academic-honesty}

The University of South Alabama's policy regarding Student Academic
Conduct Policy is found in The Lowdown
\url{http://www.southalabama.edu/lowdown/academicconductpolicy.shtml}:
The University of South Alabama is a community of scholars in which the
ideals of freedom of inquiry, freedom of thought, freedom of expression,
and freedom of the individual are sustained. The University is committed
to supporting the exercise of any right guaranteed to individuals by the
Constitution and the Code of Alabama and to educating students relative
to their responsibilities.

Violation of academic conduct policy may result in receiving 0 credit
for the affected exam/assignment.

Do not pass others work off as your own. This constitutes plagiarism and
seriously undermines your education. Students may learn about the
meaning of plagiarism and how to avoid it at the following link:
\url{http://www.southalabama.edu/univlib/instruction/plagiarismforstudents.html}

\subsection{Generative Artificial Intelligence (AI)
Policy}\label{generative-artificial-intelligence-ai-policy}

The use of Generative Artificial Intelligence (GenAI) tools (e.g.,
ChatGPT, Gemini, etc.) in this course is permitted only as explicitly
outlined below and for specific assignments. Assignments and other
graded assessments will include clear allowances for GenAI use when
assigned. Students are responsible for understanding and adhering to
these guidelines. Any use of GenAI beyond what is permitted, or without
proper acknowledgment and disclosure when usage is allowed, will be
considered a violation of the University's Student Academic Conduct
Policy and will be processed in the same manner as other academic
misconduct violations. When in doubt, always err on the side of caution
and consult with the instructor before using GenAI on graded
assignments.

Examples of generally permitted AI use:

\begin{enumerate}
\def\labelenumi{\arabic{enumi}.}
\item
  Brainstorming and building review materials
\item
  Proofing written work and checking grammar
\item
  Summarizing lecture transcripts
\item
  Editing computer code for data analysis and visualization
\end{enumerate}

Instances of unauthorized GenAI use in coursework will be treated as a
violation of the University's Student Academic Conduct Policy and will
be processed in the same manner as other academic misconduct violations
using the Academic Misconduct Penalty Record (AMPR).

\subsection{Course and Teacher
evaluation}\label{course-and-teacher-evaluation}

Student input for the purpose of course improvement is taken very
seriously and will potentially be done periodically. Please take the
time to evaluate this course and the instructor, especially at the end
of the semester. Evaluations will in no way affect your grade.

\section{Safety}\label{safety}

\subsection{Field Safety}\label{field-safety}

\subsubsection{Transportation}\label{transportation}

Wear your seatbelts at all times while in moving vehicles. Do not
distract the driver.

\subsubsection{Location awareness}\label{location-awareness}

Be aware of potential dangers (animals, weather, etc) and avoid them.
When in a group, please think about the safety of your peers.

\subsection{Building Safety}\label{building-safety}

\subsubsection{Fire}\label{fire}

Random fire safety drills may occur during the course of the term. You
will be expected to evacuate the building and assemble at the designated
location: outside by the northeast corner of the building (follow your
instructor).

\subsubsection{Tornado}\label{tornado}

Move away from external windows to the hallway in the basement, first,
or second floor.

\subsubsection{Active Shooter Incident}\label{active-shooter-incident}

``Get Out'' of the building as long as they can do so safely. Once you
have exited the building and are at a safe location, ``Call Out'' to law
enforcement using 911, USA Police at 460-6312, or press the Emergency
button on your LiveSafe Mobile Safety App. If you can't get out,
``Shelter in Place'' in a secure location that is away from windows and
doors, preferably a location that can be locked and barricaded.

\section{Campus resources}\label{campus-resources}

\subsection{Sick leave policy}\label{sick-leave-policy}

It is our responsibility to one another to avoid passing contagious
diseases as much as possible. If you are feeling ill, consider staying
home and attending remotely if possible. For lectures, you do not need
to contact me and can attend via zoom or watch the recorded lectures by
request.

If you are dealing with longer-term illness, requiring more than missing
a single day of class, please contact me to discuss accommodations.

\subsection{Writing center}\label{writing-center}

Over the course of this semester, seek help at the Writing Center if you
would like additional guidance and critique. You can access the Writing
Center at
\href{https://www.southalabama.edu/departments/academicsuccess/cae/how.it.works.html}{at
their website}. During graduate school, I found workshops and writing
groups at the Writing Center on campus to be helpful in building my
motivation and teaching me new strategies for writing efficiently.

\subsection{Disabilities}\label{disabilities}

In accordance with the Americans with Disabilities Act, students with
bona fide disabilities will be afforded reasonable accommodations. The
Office of Special Student Services (OSSS) will certify a disability and
advise faculty members of reasonable accommodations. If you have a
specific disability that qualifies you for academic accommodations,
please notify the instructor/professor and provide certification from
the Office of Special Student Services. OSSS is located at 5828 Old
Shell Road at Jaguar Drive, (251-460-7212).

\subsection{Title IX -- Sex discrimination
resources}\label{title-ix-sex-discrimination-resources}

Discrimination and harassment have no place in our community. The Title
IX office on campus is the formal location where regulations of sexual
misconduct are heard and investigated. Misconduct includes dating and
domestic violence, electronic or in person stalking, retaliation,
unequal treatment based on sex, and sexual assault or harassment.
University employees are required to report sexual misconduct to the
Title IX office that we observe or otherwise learn about. More
information on the Title IX office can be found
\href{https://www.southalabama.edu/departments/studentaffairs/titlenine/}{at
their website}.

\subsection{Food insecurity}\label{food-insecurity}

If you are subject to food insecurity, JagPantry provides a food for
students needing assistance. More information can be found
\href{https://www.southalabama.edu/departments/sga/foodpantry.html}{on
their website}

\subsection{Lower Internet Costs}\label{lower-internet-costs}

Because of the hybrid delivery mode of the course, reliable internet is
important to maintain. There are several avenues that can help provide
free or low-cost home internet some can be found through
\href{https://adeca.alabama.gov/affordable-connectivity-program/}{the
Affordable Connectivity Program} run by the State of Alabama. One
service supported through these programs is
\href{https://www.truconnect.com/findhelp}{TruConnect}. If you are
supported by a Pell Grant, are a veteran, or have other qualifying
statuses, this option can provide free broadband.

\subsection{Mental health}\label{mental-health}

Being a university student is stressful and can negatively impact mental
health. The University Counseling and Testing Center provides
confidential, free counseling and crisis intervention services to
eligible USA students and consultation and outreach services to members
of the USA community. More information can be found
\href{https://www.southalabama.edu/departments/counseling/}{at their
website}.

\subsection{Disaster plan}\label{disaster-plan}

In case of Hurricane we will follow university recommendations and
off-campus accommodations will be made. Our broad goals for the semester
will continue but mode of delivery and content may change.

In case of a shift to remote learning, we will heavily rely on Canvas.
We will have at least one synchronous lecture meeting that will be
recorded and made available through Canvas. We will use digital options
for rock samples and may also use \href{https://rockd.org/}{Rockd.org}
for local outcrop locations if we are dispersed.

\section{Final thoughts}\label{final-thoughts}

This document is a roadmap for our semester. We learn about the Earth
together and our individual experiences shape how we interpret and value
data. Like all your classes, you will get out what you put into this
course. Asking for help from one another and your instructors is
important, don't be afraid to ask a question about something you don't
know or if you want to check your knowledge about something you think
you know.

\textbf{If this document is updated, a copy will be supplied to you via
Canvas and changes will be announced in class.}




\end{document}

\makeatletter
\def\@maketitle{%
  \newpage
%  \null
%  \vskip 2em%
%  \begin{center}%
  \let \footnote \thanks
    {\fontsize{18}{20}\selectfont\raggedright  \setlength{\parindent}{0pt} \@title \par}%
}
%\fi
\makeatother